\section{Infinite volume and field calculus}\label{sec:partII-F}

We pass from the fixed-volume kernels of Theorem~\ref{prop:UV} (fixed-volume limit) to correlations on~$\R^4$. The order of limits is that of Definition~\ref{def:D09} and Figure~\ref{fig:two-limits}: at every fixed admissible volume the spacing tends to zero along the full cutoff sequence; only then does the volume grow, along one refinement of the prescribed volume sequence that serves every field list, derivative and test. Section~\ref{sec:partII-F.1} extracts this volume limit, with all derivatives, from the uniform bounds of Proposition~\ref{prop:UB} (Lemma~\ref{lem:F1}). Section~\ref{sec:partII-F.2} bounds separated correlations cell by cell and shows that negligible errors vanish in them (Lemma~\ref{lem:F1b}). Section~\ref{sec:partII-F.3} identifies the quotient fields: it constructs coefficient maps for inclusions and positional derivatives and proves $\dim_\C F_d=\dim_\C\maxwellspace{d}$ (Lemma~\ref{lem:F2}, which proves Proposition~\ref{prop:FC}). Section~\ref{sec:partII-F.4} passes reflection positivity to the limit and proves the power bounds of E7 (Lemma~\ref{lem:F3}, which completes Theorem~\ref{prop:IV}). Section~\ref{sec:partII-F.5} bounds repeated reflected insertions, the input for the slab operators of Lemma~\ref{lem:H1.1}.

The Wilson family and separated tests are those of Section~\ref{sec:continuum}, the insertion labels those of Section~\ref{sec:local-fields}, and the differences those of Section~\ref{sec:uv-limit}. Write $M_{p,v}=M^{\rm UV}_{p,v}$ and $\langle F\rangle_{N,p}=\int F\,d\mu_{N,p}$. Constants may depend on the fixed construction; further dependence is indicated.

\subsection{Volume limits}\label{sec:partII-F.1}

Fix a strictly increasing~$\mathsf p:\N\to\N$. Proposition~\ref{prop:UB} bounds the lattice kernels and their differences on separated sets, uniformly in the volume; one further difference makes these bounds equicontinuous.
\begin{lemma}[Volume limits]\label{lem:F1}
There are a strictly increasing~$\mathsf q$ and a complex multilinear table~$S$ such that, with $P=\mathsf p\circ\mathsf q$, for every~$n$, label list~$v$ and difference list~$\boldsymbol l$,
\begin{equation}\label{eq:F-conf}
S_n(v;\cdot)\in C^\infty(\Conf{n}),\qquad
\partial_{\boldsymbol l}M_{P(k),v}\longrightarrow
\partial_{\boldsymbol l}S_n(v;\cdot)
\quad\text{locally uniformly on }\Conf{n}\,.
\end{equation}
Every compact $K\subset\Conf{n}$ eventually lies in~$\NoHit{L_{P(k)}}{n}$, and for every separated test~$f$,
\begin{equation}\label{eq:F-ordered}
C_{P(k)}^{\rm UV}(v;f)\longrightarrow\int f(x)S_n(v;x)\,dx.
\end{equation}
Both convergences persist along every further refinement of~$\mathsf q$.
\end{lemma}
\begin{proof}
Fix compact sets $K\subset K'\subset\Conf{n}$, with $K'$ a neighbourhood of~$K$. For large~$p$, $K'$ lies in~$\NoHit{L_p}{n}$ with torus separation at least some $2r>0$ independent of~$p$. Proposition~\ref{prop:UB} bounds $\Delta_{\boldsymbol l}\latmoment_{N,p}$ by $B(v,\boldsymbol l,r)$ near~$K'$ for large~$N$, uniformly in~$p$, and by Lemma~\ref{lem:E3.2} the multilinear interpolants converge at each fixed~$p$ to $\partial_{\boldsymbol l}M_{p,v}$. Hence
\begin{equation}\label{eq:F-local}
\sup_{p\ge p_0}\sup_{x\in K'}|\partial_{\boldsymbol l}M_{p,v}(x)|
\le B(v,\boldsymbol l,r),
\end{equation}
although the cutoff threshold may depend on~$p$. One more derivative gives equicontinuity on~$K$. Arzelà--Ascoli~\cite[Theorem~7.25]{Rudin1976} and a diagonal extraction over compact exhaustions of~$\Conf{n}$, arities, degree lists, basis coefficient lists and difference lists, a countable family, give~$\mathsf q$. Multilinearity extends the limits to every label list, and integration by parts against compactly supported tests identifies the limit of each derivative with the derivative of the limit. This proves~\eqref{eq:F-conf}.

A compact separated set has bounded diameter, so it eventually lies in~$\NoHit{L_{P(k)}}{n}$. Theorem~\ref{prop:UV} gives, at every admissible~$p$, the full inner limit $C_p^{\rm UV}(v;f)=\int f\,M_{p,v}$; the support of~$f$ is eventually admissible along~$P$, and local uniform convergence proves~\eqref{eq:F-ordered}. Both properties pass to subsequences.
\end{proof}
Lemma~\ref{lem:F2} refines~$\mathsf q$ once more without changing~$S$; from then on~$P$ denotes the final refinement. At each volume the kernel limits remain the full ultraviolet limits of Theorem~\ref{prop:UV}: no spacing subsequence enters~\eqref{eq:F-ordered}.

\paragraph{Symmetries of the limit.} Permutation symmetry, the empty correlation, deletion of the distinguished unit, conjugation and signed-permutation covariance hold on the lattice once~$N$ exceeds the reference levels (Definition~\ref{def:D08}, Lemma~\ref{lem:E2.1a}). Approximating separated configurations by lattice points, Lemma~\ref{lem:E3.2} and~\eqref{eq:F-conf} pass them to~$S$. For translations, approximate the translation vector on the same lattice; translation invariance and~\eqref{eq:F-local} give arbitrary real translations. This proves E1--E3, and~\eqref{eq:F-conf} gives E6.

\subsection{Separated cells and negligible errors}\label{sec:partII-F.2}

Chessboard seminorms, interior observables, signed images, aligned tilings, anchors and margins are those of Definition~\ref{def:D27} and of the paragraphs before Lemma~\ref{lem:F1a} and Definition~\ref{def:E2-small}. For observables $F_1,\ldots,F_n$ in distinct cells $C_1,\ldots,C_n$ of a translated bond tiling~$\mathcal T$, each interior to its cell, write $\|F_h\|_{\mathcal T}=\|G_h\|_{\mathcal T}$, where $G_h$ is the base-cell observable with $(G_h)_{C_h}=F_h$. The chessboard inequality~\eqref{eq:E-chess}, with the constant one in the remaining cells, gives
\begin{equation}\label{eq:F-chess}
\left|\left\langle\prod_{h=1}^nF_h\right\rangle_{N,p}\right|
\le\prod_{h=1}^n\|F_h\|_{\mathcal T}\,.
\end{equation}
The seminorm measures the cost of repeating one observation throughout the box, so this bound is independent of the volume and of the distances between the cells.

Separated anchors must therefore be placed in distinct cells, away from the boundaries. In each axis, the translations of a tiling of side~$a$ that bring one of~$n$ anchors within distance~$w$ of a cell boundary form, modulo~$a$, at most~$n$ intervals of total length~$2nw$. If $2nw<a$, some lattice translation avoids them once~$a_N$ is small. We call this the strip count. Anchors at torus sup-distance greater than~$a$ never share a cell.
\begin{lemma}[Separated cells]\label{lem:F-cells}
Let $m\le k$ and $n\ge1$ with $6na_k<a_m$, and suppose the insertions~$O_{N,p}(v_h,\cdot)$, $h\le n$, satisfy the hypothesis of Lemma~\ref{lem:F1a} at level~$k$ with a common bound~$B$. For all sufficiently large~$N$ and all anchors with pairwise torus sup-distance greater than~$a_m$,
\begin{equation*}
\left|\left\langle\prod_{h=1}^nO_{N,p}(v_h,y_h)\right\rangle_{N,p}\right|\le B^n.
\end{equation*}
\end{lemma}
\begin{proof}
The strip count with $a=a_m$ and $w=3a_k$ gives a tiling of side~$a_m$ with the anchors in distinct cells, at distance at least~$3a_k$ from their boundaries. Lemma~\ref{lem:F1a} bounds each seminorm by~$B$, and~\eqref{eq:F-chess} gives the claim.
\end{proof}

A negligible error has a tolerance chosen before the cell scale (Definition~\ref{def:E2-small}), and this order makes it vanish against each fixed list of spectators.
\begin{lemma}[Vanishing errors]\label{lem:F1b}
For every negligible error~$E$ of Definition~\ref{def:E2-small}, fixed label list~$v$, distinguished slot~$i$ and~$\delta>0$, there is $C(v,\delta)<\infty$ such that, eventually in~\eqref{eq:E2-tails},
\begin{equation}\label{eq:F-small}
\left|\left\langle E_{N,p}(y_i)\prod_{h\ne i}O_{N,p}(v_h,y_h)\right\rangle_{N,p}\right|
\le C(v,\delta)\eta_{N,p}
\qquad(\operatorname{sep}_{\rm tor}(y)\ge\delta)\,.
\end{equation}
Consequently its smear against any fixed separated test tends to zero in that order.
\end{lemma}
The error and the spectators receive distinct cells; the error pays its tolerance~$\eta_{N,p}$, and each spectator a constant.
\begin{proof}
Let~$n$ be the arity. Take~$m$ with $a_m<\delta$, and let $k\ge m$ be the level that Definition~\ref{def:E2-small} assigns to the arity~$n$ and the lower level~$m$; the error tiling has side~$a_k$. Choose $k_{\rm sp}\ge k$, beyond the reference levels of the spectators, with $3a_{k_{\rm sp}}\le a_k/[8(n+1)]$. Lemma~\ref{lem:E2.3} bounds the centered seminorms of all signed and conjugated images of the spectators at level~$k_{\rm sp}$ by some~$B_{\rm sp}$. The strip count with $a=a_k$ and $w=a_k/[4(n+1)]$, and $\delta>a_k$, give a tiling of side~$a_k$ with the anchors in distinct cells at distance at least~$w$ from their boundaries. This margin contains the integer margin~\eqref{eq:E2-margin}, including its extra lattice unit, and exceeds~$3a_{k_{\rm sp}}$. Translating by a multiple of~$a_k$, we may take the error's cell to be the base cell. Definition~\ref{def:E2-small} bounds the seminorm of the error by~$\eta_{N,p}$, and Lemma~\ref{lem:F1a}, with tiling side~$a_k$ and centered level~$k_{\rm sp}$, bounds those of the spectators by~$B_{\rm sp}$. Now~\eqref{eq:F-chess} gives~\eqref{eq:F-small} with $C=B_{\rm sp}^{n-1}$. On a compact test support the number of lattice tuples times~$a_N^{4n}$ is bounded independently of~$N,p$, which proves the smeared assertion.
\end{proof}

\subsection{Quotient fields}\label{sec:partII-F.3}

Two labels represent the same field when every separated correlation gives the same answer for them (Definition~\ref{def:D10}). For a degree~$d$, a slot, fixed spectators~$v$ and a separated test~$f$, the one-slot functionals on~$V_d$ are
\begin{equation*}
L_{N,p}(c)=\latsmear_{N,p}\bigl((d,c),v;f\bigr),\qquad
L(c)=\int f(x_0,x)\,S_{n+1}\bigl(c(x_0),v(x)\bigr)\,dx_0\,dx.
\end{equation*}
As $S$ is smooth, $\nullspace{d}(S)$ is the intersection of the kernels of all such~$L$: it consists exactly of the labels that no separated smear detects. Four inputs from Section~\ref{sec:partII-E.2} identify it, each with one task. (C1)~Lemma~\ref{lem:E2.8} represents a degree-$d$ insertion in a degree $e\ge d$, and (C2)~Lemma~\ref{lem:E2.9} its lattice difference quotient in degree~$d+1$, up to negligible errors. (C3)~Lemma~\ref{lem:E2.7} replaces a tame family in $\ker\Pi_{d,N,p}$ by observables with vanishing separated moments, up to a negligible error; \emph{tame} means the subpower bound~\eqref{eq:F-subpower}. (C4)~The free test rows of Lemma~\ref{lem:E2-free-tests} and the relative rank comparison of Lemma~\ref{lem:E2.10} keep the quotient from losing dimension. The families supplied by (C1)--(C3) may have unbounded coefficients; Lemma~\ref{lem:F-represent} turns them into fixed coefficient vectors.

The coefficient limits are taken along a common spacing subsequence~$\mathsf n$ at every volume and a further volume refinement, in the order~\eqref{eq:E2-tails}. They do not affect the kernel limits, which at each admissible volume remain the full ultraviolet limits required by Definition~\ref{def:D09}. A family in a finite-dimensional space is \emph{bounded} if its norm is bounded eventually in~\eqref{eq:E2-tails}; an \emph{ordered limit point} of countably many bounded families is their joint limit along an ultrafilter containing the ordered tails. It exists by compactness, and families that converge along the ordered tails keep their limits, so passing to ordered limit points changes neither~$S$ nor~$P$.

The next lemma, finite-dimensional linear algebra, converts approximate representatives, which need not be bounded, into one exact representative.
\begin{lemma}[Representation of limits]\label{lem:F-represent}
Let~$V$ be a finite-dimensional complex normed space and~$\mathcal F$ a proper filter on an index set. Let $\Pi_u,\Pi$ be idempotents on~$V$ with $\Pi_u\to\Pi$, and $L_{i,u},L_i$, $i\in I$, linear functionals with $L_{i,u}\to L_i$ and $\ker\Pi=\bigcap_i\ker L_i$. Suppose $x_u\in V$ satisfy, for every~$i$, along~$\mathcal F$,
\begin{equation}\label{eq:F-normalized}
L_{i,u}(x_u)\longrightarrow t_i,\qquad
L_{i,u}\bigl((x_u-\Pi_ux_u)/(1+\|\Pi_ux_u\|)\bigr)\longrightarrow0.
\end{equation}
Then some $z\in V$ satisfies $\Pi z=z$ and $L_i(z)=t_i$ for every~$i$.
\end{lemma}
The second condition makes the null part of~$x_u$ negligible relative to its retained part. If the retained parts~$\Pi_ux_u$ escaped to infinity, normalizing them would produce a unit vector in the range of~$\Pi$ annihilated by every~$L_i$: a null vector in the retained range, which $\ker\Pi=\bigcap_i\ker L_i$ excludes.
\begin{proof}
Put $b_u=(1+\|\Pi_ux_u\|)^{-1}$ and $z_u=b_u\Pi_ux_u$, so $b_u+\|z_u\|=1$ and $\Pi_uz_u=z_u$. The pairs $(b_u,z_u)$ have a cluster point~$(\lambda,z)$ along~$\mathcal F$ in $[0,1]\times\{\|z\|\le1\}$. As $L_{i,u}(z_u)=b_uL_{i,u}(x_u)-L_{i,u}(b_u(x_u-\Pi_ux_u))$,
\begin{equation*}
L_i(z_u)-b_ut_i=b_u\bigl(L_{i,u}(x_u)-t_i\bigr)-L_{i,u}\bigl(b_u(x_u-\Pi_ux_u)\bigr)+(L_i-L_{i,u})(z_u)\longrightarrow0,
\end{equation*}
and $\Pi z_u-z_u=(\Pi-\Pi_u)z_u\to0$. Hence $L_i(z)=\lambda t_i$, $\Pi z=z$ and $\lambda+\|z\|=1$. If $\lambda=0$, then $z\in\bigcap_i\ker L_i=\ker\Pi$, so $z=\Pi z=0$, contradicting $\|z\|=1$. Thus $\lambda>0$, and $z/\lambda$ has the required properties.
\end{proof}

\begin{lemma}[Corrected coefficient maps]\label{lem:F2}
After a strictly increasing refinement of~$\mathsf q$, which does not change~$S$, there are complex linear maps $\hat\jmath_{d,e}:V_d\to V_e$ for $d\le e$ and $\hat\partial_{d,a}:V_d\to V_{d+1}$ such that, for every~$c\in V_d$, spectator word~$v(x)$ and $(x_0,x)\in\Conf{n+1}$,
\begin{equation}\label{eq:F-calculus}
\begin{gathered}
S_{n+1}\bigl(\hat\jmath_{d,e}c(x_0),v(x)\bigr)=S_{n+1}\bigl(c(x_0),v(x)\bigr),\\
S_{n+1}\bigl(\hat\partial_{d,a}c(x_0),v(x)\bigr)=\partial_{x_0^a}S_{n+1}\bigl(c(x_0),v(x)\bigr).
\end{gathered}
\end{equation}
Moreover $\dim_\C F_d=\dim_\C\maxwellspace{d}$, and the induced maps $\jmath_{d,e}[c]=[\hat\jmath_{d,e}c]$ and $\partial_a[c]=[\hat\partial_{d,a}c]$ satisfy Proposition~\ref{prop:FC}.
\end{lemma}
The proof identifies the null space with the kernel of a limiting projection, then applies Lemma~\ref{lem:F-represent} to the families of (C1) and~(C2).
\begin{proof}
\emph{Limiting projections.} Bounded coefficient families are tame: Lemma~\ref{lem:E2.2} bounds their inverse terminal coordinates by a fixed factor times $\prod_{m<N}C^{\rm col}_{J_{m+1}}$, and $\log C^{\rm col}_{J_{m+1}}\le J_{m+1}^{-2}\log\ell_m$ (Appendix~\ref{sec:constants}) with $J_m\to\infty$ makes this product grow more slowly than every power of~$a_N^{-1}$. A countable diagonal extraction, first in~$N$ at every~$p$, then in the volume, retaining~\eqref{eq:F-conf}, makes the following converge: the projections of Lemma~\ref{lem:E2.6}, which are bounded by~$C^\Pi_d$, the comparison lifts of Lemma~\ref{lem:E2.10}, and the couplings $\beta_{N,p,k}\in\mathfrak B_k$, $k\in\N$, of Theorem~\ref{prop:MT}. In operator norm,
\begin{equation}\label{eq:F-projections}
\Pi_{d,\mathsf n(j),p}\longrightarrow\Pi_d^{(p)},\qquad
\Pi_d^{(P(k))}\longrightarrow\Pi_d,\qquad
\Pi_d^2=\Pi_d\,,
\end{equation}
and by~\eqref{eq:E2-tails} this iterated limit is an ordered limit.

\emph{The null space.} By finite dimensionality, the full fixed-volume limits and~\eqref{eq:F-ordered}, $L_{N,p}\to L$ in operator norm in the ordered limit. For fixed~$c$, the bounded family $c-\Pi_{d,N,p}c$ is tame and lies in $\ker\Pi_{d,N,p}$, so (C3) and Lemma~\ref{lem:F1b} give $L_{N,p}(c-\Pi_{d,N,p}c)\to0$. Thus $L((1-\Pi_d)c)=0$ for every~$L$:
\begin{equation}\label{eq:F-null}
\ker\Pi_d\subset\nullspace{d}(S).
\end{equation}
For the reverse inequality, use (C4). Rearranging~\eqref{eq:F-relative} isolates~$\mathcal G^\nu_I$; every other term contains a nonempty~$R_J$, whose expectation is at most its seminorm~$\gamma_\nu$ by~\eqref{eq:F-chess} with one occupied cell. Thus
\begin{equation}\label{eq:F-rank}
\left|\left\langle\prod_{h\in I}F^{\rm tst}_h\right\rangle_{N,p}-\mathcal G^\nu_I\right|
\le\gamma_\nu\sum_{J\subsetneq I}|\mathcal G^\nu_J|,
\end{equation}
and the sum stays bounded as $y^{\nu,i}\to y^i$. Pass to the coefficient and kernel limits. Each entry of the resulting matrix is a fixed spectator functional, vanishing on~$\nullspace{d}(S)$, applied to the limiting lift of~$\mathfrak m_j$; so the matrix factors through~$F_d$ and has rank at most~$\dim F_d$. As $\nu\to\infty$ these matrices tend to the rank-$\dim_\C\maxwellspace{d}$ matrix~$\mathcal G$, so $\dim F_d\ge\dim_\C\maxwellspace{d}$. Ranks of converging idempotents are eventually constant (if $\|P-Q\|<1$, then $P$ is injective on $\operatorname{ran}Q$), so the dimension inequality of Lemma~\ref{lem:E2.6} passes to~$\Pi_d$, and with~\eqref{eq:F-null}, $\dim F_d\le\dim V_d-\dim\ker\Pi_d\le\dim_\C\maxwellspace{d}$. Hence
\begin{equation}\label{eq:F-dimension}
\nullspace{d}(S)=\ker\Pi_d,\qquad \dim_\C F_d=\dim_\C\maxwellspace{d}.
\end{equation}

\emph{Targets.} For the inclusion of $c\in V_d$, (C1) and Lemma~\ref{lem:F1b} give a tame family $w_{N,p}\in V_e$ whose one-slot functionals tend to $t_L=\int fS(c,\ldots)$. For the derivative in slot~$i$ and axis~$a$, (C2) gives $w_{N,p}\in V_{d+1}$, and torus summation by parts gives exactly
\begin{equation}\label{eq:F-parts}
a_N^{4n}\sum_y f(\hat y)\Delta_{i,a}\latmoment_{N,p}(v;y)
=a_N^{4n}\sum_y\frac{f(\hat y-a_Ne_{i,a})-f(\hat y)}{a_N}\latmoment_{N,p}(v;y),
\end{equation}
with $e_{i,a}$ the coordinate vector of slot~$i$, axis~$a$. The backward quotient tends to~$-\partial_{i,a}f$ with residual at most $a_N\|\partial_{i,a}^2f\|_\infty/2$, which Proposition~\ref{prop:UB} at order zero controls, so the target is $t_L=-\int(\partial_{i,a}f)S(c,\ldots)$.

\emph{Representation.} Apply Lemma~\ref{lem:F-represent} in~$V_e$, with~$\mathcal F$ the filter of ordered tails $\{(j,k):k\ge k_0,\ j\ge j_0(k)\}$, $\Pi_u=\Pi_{e,\mathsf n(j),P(k)}$ and the one-slot functionals of degree~$e$. They converge as shown above, and $\ker\Pi_e=\bigcap\ker L$ by~\eqref{eq:F-dimension}. With $\Pi=\Pi_{e,N,p}$, the family $(w-\Pi w)/(1+\|\Pi w\|)$ lies in $\ker\Pi$ and is tame, being a bounded multiple of~$w$ minus a family of norm at most one; (C3) and Lemma~\ref{lem:F1b} give the second condition of~\eqref{eq:F-normalized}. The lemma gives one vector representing the target against every spectator and test at once; no further subsequence is chosen for these uncountably many tests. Choose such vectors on a basis of~$V_d$ and extend linearly; \eqref{eq:F-conf} turns the distributional identities into~\eqref{eq:F-calculus}.

\emph{Induced maps.} By~\eqref{eq:F-calculus}, $\hat\jmath_{d,e}$ and $\hat\partial_{d,a}$ carry $\nullspace{d}(S)$ into null spaces, so they descend to~$F_d$ and satisfy the identities of Proposition~\ref{prop:FC}; with~\eqref{eq:F-dimension} this proves that proposition. The further properties listed after it (injective responses, units, conjugation, $\hypercubic$ covariance, the positional identity in every slot) follow from these identities, E1 and E2, as explained there.
\end{proof}

\paragraph{Example: the curvature square (continued).} Section~\ref{sec:A-example} follows the label $c_{F^2}\in V_4$ of $\frac14\sum_{a,b,\alpha}F^\alpha_{ab}F^\alpha_{ab}$ to its insertion; in degree four the quotient removes nothing. Every jet coordinate has weight at least two, and since~$\g$ is simple no nonzero $\operatorname{Ad}(G)$-invariant polynomial is linear in a $\g$-valued jet; so an element of~$V_4$ is a constant plus an invariant quadratic form in the lowest jet~$\mathcal A_0$, that is, in the curvature at the origin. The Maxwell equations do not constrain the curvature at one point, so no nonzero such polynomial vanishes on every Maxwell solution: the equation space~$\mathcal I_4$ of Lemma~\ref{lem:E2.6} is zero. That lemma gives $\Pi_{4,N,p}=1$, and~\eqref{eq:F-dimension} gives $\nullspace{4}(S)=0$. Distinct degree-four labels give distinct fields, and $[c_{F^2}]\ne0$ in~$F_4$. The quotient first acts on derivatives: the lattice difference quotient of the insertion of~$c_{F^2}$ is not a degree-five insertion, (C2) represents it by a tame family in~$V_5$, and Lemma~\ref{lem:F-represent} replaces that family by one vector, with $\partial_a[c_{F^2}]=[\hat\partial_{4,a}c_{F^2}]$. In degree five, $\mathcal I_5$ contains the invariant contraction $\sum_{a,b,c,\alpha}F^\alpha_{ab}\,\partial_cF^\alpha_{cb}$ of a curvature with the Maxwell expression. By Lemma~\ref{lem:E2.6} and~\eqref{eq:F-dimension}, $\nullspace{5}(S)=\ker\Pi_5$ is therefore nonzero, and $\hat\partial_{4,a}c_{F^2}$ is determined only modulo these limits of transported equation directions.

\subsection{Euclidean properties}\label{sec:partII-F.4}

Lattice symmetries and positivity pass through the ordered limits; Lemma~\ref{lem:F-cells} controls separated points as they approach one another, uniformly in their diameter.
\begin{lemma}[Euclidean properties]\label{lem:F3}
The table~$S$ of Lemma~\ref{lem:F1} satisfies E1--E7.
\end{lemma}
\begin{proof}
E1--E3 and E6 were proved after Lemma~\ref{lem:F1}; the positive derivative orders of E7 use~\eqref{eq:F-calculus}.

\emph{Reflection positivity.} We prove E5 and deduce E4. Fix joint smooth tests~$f_i$ with compact support on positive-side separated configurations for an axis~$a$. Use the bond reflection~$x^a\mapsto a_N-x^a$, a lattice symmetry with planes~$x^a=a_N/2$ and~$x^a=(a_N+L_p)/2$. For large~$p,N$ the lattice smears~$F_i$ are supported strictly inside the positive half-torus, with stencils avoiding both planes, since the supports keep a positive distance from $x^a=0$ and the stencils have physical radius tending to zero. Periodic Wilson reflection positivity~\cite[Theorem~2.1, pp.~14--15]{Seiler1982} gives
\begin{equation*}
\sum_{i,j}\bar z_i z_j\langle\overline{F_i\circ\mathsf r_{a,N}}F_j\rangle_{N,p}\ge0,
\end{equation*}
where~$\mathsf r_{a,N}$ reflects configurations. Insertion covariance reflects the positions and applies~$\mathsf s_a^{\rm refl}$ to the labels; conjugation acts on the labels of the first word. The reflected joint tests converge with all derivatives to~$\bar f_i\otimes f_j$ with the first group reflected through~$x^a=0$, and all supports lie in one compact separated set. Proposition~\ref{prop:UB} at order zero bounds the effect of the moving tests by their uniform difference times a bounded weighted site count. Lemma~\ref{lem:E3.2} and~\eqref{eq:F-ordered} then give a positive semidefinite limiting matrix; smoothness of~$S$ makes every pairing absolutely integrable. This is E5. Concentrating normalized bumps at positive-side words gives E4 by continuity.

\emph{Power bound at order zero.} Take a nonempty label list and put $d_{\rm pow}=1+\max_id_i$. Lemma~\ref{lem:E2.3} with $\varepsilon=1$, for the finitely many signed and conjugated images of the labels, gives $k_0$ and $C_{\rm ctr}$ such that their centered seminorms at levels $k\ge k_0$ are at most $C_{\rm ctr}a_k^{-d_{\rm pow}}$; these constants precede the separation. The block ratios grow slowly: $\log\ell_u/s_u\to0$ (Lemma~\ref{lem:E1.1}) gives $\ell_u\le a_u^{-1}$ eventually, so for some $c_{\rm mesh}>0$ and all~$u$
\begin{equation}\label{eq:F-thickness}
a_{u+1}\ge c_{\rm mesh}a_u^2.
\end{equation}
Thus the first $u\ge k_0$ with $a_u<\delta\le1$ satisfies $a_u\ge c_0\delta^2$, where $c_0=\min(c_{\rm mesh},a_{k_0})$: a discrete scale lies within a fixed power of any prescribed length. Fix $0<r\le1$ and~$x$ with $\sepinf(x)\ge r$. Choose the first $m\ge k_0$ with $a_m<r/2$, then the first $k\ge k_0$ with $a_k<a_m/[6(n+1)]$; then $k\ge m$ and $a_k\ge c'(n)r^4$. Lattice approximations of~$x$ in large volumes have torus separation greater than~$a_m$, so Lemma~\ref{lem:F-cells} bounds the lattice correlation by $(C_{\rm ctr}a_k^{-d_{\rm pow}})^n$. Lemma~\ref{lem:E3.2} and~\eqref{eq:F-conf} give
\begin{equation}\label{eq:F-power}
|S_n(v;x)|\le C(n,d,c)r^{-4d_{\rm pow}n},
\end{equation}
with constants independent of the positions and the volume. Each insertion costs its centered seminorm at a cell scale comparable to~$r^4$, and the chessboard bound makes the total independent of the diameter of~$x$.

\emph{Higher orders.} By~\eqref{eq:F-calculus}, every coordinate derivative is a correlation with fixed degree-raised labels, finitely many at each order. Apply~\eqref{eq:F-power} to each, take the largest exponent, and bound the Fr\'echet norm by the sum of the coordinate coefficients in an orthonormal basis of~$\R^{4n}$. This proves E7; for arity zero, $S_0=1$.
\end{proof}

\begin{proof}[Proof of Theorem~\ref{prop:IV}]
Compose the volume refinements of Lemmas~\ref{lem:F1} and~\ref{lem:F2}. Equation~\eqref{eq:F-ordered} proves ordered convergence, with the full inner limit of Theorem~\ref{prop:UV} at every admissible volume, and Lemma~\ref{lem:F3} proves E1--E7. The coefficient corrections change no kernel limit.
\end{proof}

\subsection{Repeated reflected words}\label{sec:partII-F.5}

Use the OS space of Definition~\ref{def:D14}. For a word~$w$, let~$w^*$ reverse it and apply conjugation to its labels, and let $\theta w=\mathsf s_4^{\rm refl}w^*$ be its antilinear OS reflection. For $s>0$ and a word~$w$ with times in $(0,s)$ set
\begin{equation*}
A_wv=w\,\tau_{se_4}v,\qquad
w^\dagger=\tau_{se_4}\theta w;
\end{equation*}
$A_w$ acts on words, and boundedness on OS classes is not assumed. A word has buffer~$r$ in a slab of width~$h_{\rm sl}$ if its points have pairwise sup-distance at least~$r$ and times in $[r,h_{\rm sl}-r]$. Fix a degree bound~$d$ and a coefficient bound~$R\ge0$, and put $Q=4(d+1)$. The next estimate bounds repeated insertion before any operator extension.
\begin{lemma}[Repeated reflected words]\label{lem:F4}
There is $C_0=C_0(d,R)>0$ such that every nonempty length-$n$ word~$w$ with these bounds and buffer~$0<r\le1$ in a slab of width~$h_{\rm sl}$ admits a width $0<s\le h_{\rm sl}$ in which~$w$ has buffer~$r/2$ and, for every separated exterior word~$v$ with positive times,
\begin{equation*}
\|[(A_{w^\dagger}A_w)^jv]\|_{\rm OS}^2
\le B_v[C_0^n(n/r)^{Qn}]^{4j}\qquad(j\ge0),
\end{equation*}
where $B_v=B_v(w,r,s)<\infty$ is independent of~$j$.
\end{lemma}
The operators use the chosen width~$s$, and~$s$, $C_0$ and~$Q$ precede~$v$. The main insertions occupy cells of a size chosen from~$r$ alone and the exterior insertions a finer aligned tiling, so the cost of the exterior does not grow with the number~$j$ of repetitions.
\begin{proof}
Enlarge the coefficient ball to contain all signed and conjugated images in degrees at most~$d$. Lemma~\ref{lem:E2.3} with $\varepsilon=1$ gives $k_*$ and $B_*>0$ such that their centered seminorms at levels $k_{\rm w}\ge k_*$ are at most $B_* a_{k_{\rm w}}^{-d_{\rm pow}}$, $d_{\rm pow}=d+1$; this precedes $n,r,w$. Choose the first $m\ge k_*$ with $a_m<r/32$ and put $s=a_m\lfloor h_{\rm sl}/a_m\rfloor$. As~$w$ is nonempty with buffer~$r$, $h_{\rm sl}\ge2r$, so $0<s\le h_{\rm sl}$ and $h_{\rm sl}-s<a_m$, and~$w$ has buffer~$r/2$ in width~$s$. By~\eqref{eq:F-thickness}, $a_m\ge c_1r^2$ with $c_1>0$ depending on~$k_*$. Fix $0<\alpha\le1$ with $20\alpha<c_1$, and choose the first $k_{\rm w}\ge k_*$ with $a_{k_{\rm w}}<\alpha(r/n)^2$. Then
\begin{equation}\label{eq:F-repeat-scales}
20na_{k_{\rm w}}<a_m,\qquad a_{k_{\rm w}}\ge c_2(r/n)^4,\qquad
(B_*a_{k_{\rm w}}^{-d_{\rm pow}})^n\le (B_*c_2^{-d_{\rm pow}})^n(n/r)^{4d_{\rm pow}n}.
\end{equation}

Expanding the OS norm gives a correlation of $4jn$ main insertions and two copies of~$v$. The main insertions are separated from one another and from the exterior by at least $r/2>a_m$. Translations by multiples of~$s$ preserve residues modulo~$a_m\Z^4$, so the main anchors have at most~$2n$ residue patterns, those of~$w$ and~$\theta w$, for every~$j$. For nonempty~$v$, let $t_{\min}>0$ be its minimum time and choose $m_{\rm v}\ge k_{\rm w}$ with $a_{m_{\rm v}}<2t_{\min}$ and $a_{m_{\rm v}}$ below every distance between distinct anchors of~$v$. The two exterior copies have times~$\pm(2js+t)$, $t$ a time of~$v$, so they are more than $2t_{\min}>a_{m_{\rm v}}$ apart for every~$j$. Choose $k_{\rm v}\ge m_{\rm v}$, beyond the exterior reference levels, with $12|v|a_{k_{\rm v}}<a_{m_{\rm v}}$, and let $X\ge0$ bound the centered seminorms of the exterior labels and their images at level~$k_{\rm v}$ (Lemma~\ref{lem:E2.3}). Neither $m_{\rm v}$ nor~$k_{\rm v}$ depends on~$j$.

The coarse tiling has side~$a_m$ and separates the main insertions; the fine tiling has side~$a_{m_{\rm v}}$, shares its origin and separates the exterior insertions. Their side ratio is odd, so coarse reflections preserve the fine grid and its interior margins. The strip count for the $2n$ main residues modulo~$a_m$ with $w=5a_{k_{\rm w}}$, possible by~\eqref{eq:F-repeat-scales}, keeps the main anchors at distance at least~$5a_{k_{\rm w}}$ from the coarse boundaries. A further translation by less than~$a_{m_{\rm v}}\le a_{k_{\rm w}}$ puts the $2|v|$ exterior anchors at distance at least~$3a_{k_{\rm v}}$ from the fine boundaries, since $12|v|a_{k_{\rm v}}<a_{m_{\rm v}}$, and leaves the main anchors margin at least~$3a_{k_{\rm w}}$. The main anchors occupy distinct coarse cells, none containing an exterior anchor.

Apply~\eqref{eq:F-chess} to the coarse tiling, treating the product of the exterior factors in one coarse cell as one observable. Lemma~\ref{lem:F1a} bounds each main factor by $B_*a_{k_{\rm w}}^{-d_{\rm pow}}$. For a coarse cell with~$n_C$ exterior factors, reflect their product into every coarse cell; in each copy the exterior anchors are more than~$a_{m_{\rm v}}$ apart and occupy distinct fine cells, so~\eqref{eq:F-chess} for the fine tiling and Lemma~\ref{lem:F1a} bound the coarse seminorm of that product by~$X^{n_C}$. Therefore
\begin{equation*}
\|[(A_{w^\dagger}A_w)^jv]\|_{\rm OS}^2
\le X^{2|v|}(B_*a_{k_{\rm w}}^{-d_{\rm pow}})^{4jn}.
\end{equation*}
For each fixed~$j$, lattice approximations of the anchors preserving these strict margins, Lemma~\ref{lem:E3.2} and~\eqref{eq:F-conf} pass this bound to the limit; no cutoff threshold uniform in~$j$ is needed. Take $B_v=\max(1,X^{2|v|})$, $C_0=B_*c_2^{-d_{\rm pow}}$ and $Q=4d_{\rm pow}$, and use~\eqref{eq:F-repeat-scales}. For empty~$v$ omit its factors and take $B_v=1$.
\end{proof}
